% 图 2-3 跨节点部署 v2 — 修：auto-pin 公式移到右下角小注释 / ② 气泡远离 worker box 边缘
\documentclass[border=4pt]{standalone}
\usepackage{xeCJK}
\setCJKmainfont{Songti SC}
\setCJKsansfont{Heiti SC}
\usepackage{tikz}
\usepackage{amsmath}
\usetikzlibrary{arrows.meta, positioning, calc, shapes.geometric, fit, backgrounds}
% 共享 TikZ 风格 — Aegaeon (SOSP'25) inspired，对标顶刊系统论文
% 使用：% 共享 TikZ 风格 — Aegaeon (SOSP'25) inspired，对标顶刊系统论文
% 使用：% 共享 TikZ 风格 — Aegaeon (SOSP'25) inspired，对标顶刊系统论文
% 使用：\input{_style.tex} 在 standalone preamble 内
% 调色：Okabe-Ito（与 matplotlib 图一致）

\definecolor{fuxiBlue}{HTML}{0072B2}    % 自家系统主色
\definecolor{fuxiOrange}{HTML}{E69F00}  % 第一对照
\definecolor{fuxiGreen}{HTML}{009E73}
\definecolor{fuxiPurple}{HTML}{CC79A7}
\definecolor{fuxiSky}{HTML}{56B4E9}
\definecolor{fuxiGrayBox}{HTML}{F2F2F2} % 浅灰底
\definecolor{fuxiGrayLine}{HTML}{808080}
\definecolor{fuxiGrayDark}{HTML}{555555}

\tikzset{
  % ── 通用节点 ────────────────────────────────────────
  modBox/.style = {
    draw=fuxiGrayLine, thick, fill=white,
    rounded corners=2pt, inner sep=4pt,
    minimum height=8mm, minimum width=20mm,
    align=center, font=\small
  },
  highlightBox/.style = {
    modBox, draw=fuxiBlue, fill=fuxiBlue!8,
    text=fuxiBlue, font=\small\bfseries
  },
  ghostBox/.style = {
    modBox, draw=fuxiGrayLine!50, fill=fuxiGrayBox,
    text=fuxiGrayDark, font=\small\itshape
  },
  layerBg/.style = {
    draw=fuxiGrayLine!40, dashed, fill=fuxiGrayBox!40,
    rounded corners=4pt, inner sep=8pt
  },
  externalBoundary/.style = {
    draw=fuxiGrayLine, dashed, dash pattern=on 4pt off 3pt,
    rounded corners=6pt, inner sep=10pt
  },
  % ── 箭头 ─────────────────────────────────────────────
  flow/.style = {
    -{Stealth[length=4pt, width=4pt]},
    line width=0.7pt, color=fuxiGrayDark
  },
  flowEvent/.style = {
    flow, color=fuxiBlue, line width=0.9pt
  },
  flowDup/.style = {
    flow, color=fuxiBlue!70, line width=0.7pt, dotted
  },
  bidir/.style = {
    {Stealth[length=4pt]}-{Stealth[length=4pt]},
    line width=0.7pt, color=fuxiGrayDark
  },
  % ── 编号气泡 ① ② ③ ────────────────────────────────
  numBubble/.style = {
    circle, draw=fuxiBlue, fill=white, text=fuxiBlue,
    font=\scriptsize\bfseries, inner sep=0pt,
    minimum size=4mm, line width=0.6pt
  },
}
 在 standalone preamble 内
% 调色：Okabe-Ito（与 matplotlib 图一致）

\definecolor{fuxiBlue}{HTML}{0072B2}    % 自家系统主色
\definecolor{fuxiOrange}{HTML}{E69F00}  % 第一对照
\definecolor{fuxiGreen}{HTML}{009E73}
\definecolor{fuxiPurple}{HTML}{CC79A7}
\definecolor{fuxiSky}{HTML}{56B4E9}
\definecolor{fuxiGrayBox}{HTML}{F2F2F2} % 浅灰底
\definecolor{fuxiGrayLine}{HTML}{808080}
\definecolor{fuxiGrayDark}{HTML}{555555}

\tikzset{
  % ── 通用节点 ────────────────────────────────────────
  modBox/.style = {
    draw=fuxiGrayLine, thick, fill=white,
    rounded corners=2pt, inner sep=4pt,
    minimum height=8mm, minimum width=20mm,
    align=center, font=\small
  },
  highlightBox/.style = {
    modBox, draw=fuxiBlue, fill=fuxiBlue!8,
    text=fuxiBlue, font=\small\bfseries
  },
  ghostBox/.style = {
    modBox, draw=fuxiGrayLine!50, fill=fuxiGrayBox,
    text=fuxiGrayDark, font=\small\itshape
  },
  layerBg/.style = {
    draw=fuxiGrayLine!40, dashed, fill=fuxiGrayBox!40,
    rounded corners=4pt, inner sep=8pt
  },
  externalBoundary/.style = {
    draw=fuxiGrayLine, dashed, dash pattern=on 4pt off 3pt,
    rounded corners=6pt, inner sep=10pt
  },
  % ── 箭头 ─────────────────────────────────────────────
  flow/.style = {
    -{Stealth[length=4pt, width=4pt]},
    line width=0.7pt, color=fuxiGrayDark
  },
  flowEvent/.style = {
    flow, color=fuxiBlue, line width=0.9pt
  },
  flowDup/.style = {
    flow, color=fuxiBlue!70, line width=0.7pt, dotted
  },
  bidir/.style = {
    {Stealth[length=4pt]}-{Stealth[length=4pt]},
    line width=0.7pt, color=fuxiGrayDark
  },
  % ── 编号气泡 ① ② ③ ────────────────────────────────
  numBubble/.style = {
    circle, draw=fuxiBlue, fill=white, text=fuxiBlue,
    font=\scriptsize\bfseries, inner sep=0pt,
    minimum size=4mm, line width=0.6pt
  },
}
 在 standalone preamble 内
% 调色：Okabe-Ito（与 matplotlib 图一致）

\definecolor{fuxiBlue}{HTML}{0072B2}    % 自家系统主色
\definecolor{fuxiOrange}{HTML}{E69F00}  % 第一对照
\definecolor{fuxiGreen}{HTML}{009E73}
\definecolor{fuxiPurple}{HTML}{CC79A7}
\definecolor{fuxiSky}{HTML}{56B4E9}
\definecolor{fuxiGrayBox}{HTML}{F2F2F2} % 浅灰底
\definecolor{fuxiGrayLine}{HTML}{808080}
\definecolor{fuxiGrayDark}{HTML}{555555}

\tikzset{
  % ── 通用节点 ────────────────────────────────────────
  modBox/.style = {
    draw=fuxiGrayLine, thick, fill=white,
    rounded corners=2pt, inner sep=4pt,
    minimum height=8mm, minimum width=20mm,
    align=center, font=\small
  },
  highlightBox/.style = {
    modBox, draw=fuxiBlue, fill=fuxiBlue!8,
    text=fuxiBlue, font=\small\bfseries
  },
  ghostBox/.style = {
    modBox, draw=fuxiGrayLine!50, fill=fuxiGrayBox,
    text=fuxiGrayDark, font=\small\itshape
  },
  layerBg/.style = {
    draw=fuxiGrayLine!40, dashed, fill=fuxiGrayBox!40,
    rounded corners=4pt, inner sep=8pt
  },
  externalBoundary/.style = {
    draw=fuxiGrayLine, dashed, dash pattern=on 4pt off 3pt,
    rounded corners=6pt, inner sep=10pt
  },
  % ── 箭头 ─────────────────────────────────────────────
  flow/.style = {
    -{Stealth[length=4pt, width=4pt]},
    line width=0.7pt, color=fuxiGrayDark
  },
  flowEvent/.style = {
    flow, color=fuxiBlue, line width=0.9pt
  },
  flowDup/.style = {
    flow, color=fuxiBlue!70, line width=0.7pt, dotted
  },
  bidir/.style = {
    {Stealth[length=4pt]}-{Stealth[length=4pt]},
    line width=0.7pt, color=fuxiGrayDark
  },
  % ── 编号气泡 ① ② ③ ────────────────────────────────
  numBubble/.style = {
    circle, draw=fuxiBlue, fill=white, text=fuxiBlue,
    font=\scriptsize\bfseries, inner sep=0pt,
    minimum size=4mm, line width=0.6pt
  },
}


\begin{document}
\begin{tikzpicture}[font=\sffamily, node distance=8mm and 16mm]

% ── Home 节点（左）──
\node[highlightBox, minimum width=32mm] (home_xuannv) {玄女 + 编排器};
\node[modBox, below=of home_xuannv, minimum width=32mm] (home_dist) {dist controller\\\scriptsize axum /api + /dist};
\node[modBox, below=of home_dist, minimum width=32mm] (home_sqlite) {SQLite WAL};
\node[modBox, below=of home_sqlite, minimum width=32mm] (home_im) {IM 后端};

% ── Worker A（右上）──
\node[modBox, right=46mm of home_xuannv, minimum width=30mm] (worker1_cc) {cc 门客};
\node[modBox, below=2mm of worker1_cc, minimum width=30mm] (worker1_sandbox) {git worktree\\\scriptsize 隔离工作区};
\node[modBox, below=2mm of worker1_sandbox, minimum width=30mm] (worker1_relay) {fuxi-cli worker\\\scriptsize relay 进程};

% ── Worker B（右下）──
\node[modBox, below=18mm of worker1_relay, minimum width=30mm] (worker2_cc) {codex 门客};
\node[modBox, below=2mm of worker2_cc, minimum width=30mm] (worker2_relay) {fuxi-cli worker};

% ── 节点边界（虚线圈 + 节点标题）──
\begin{scope}[on background layer]
  \node[externalBoundary, fit=(home_xuannv) (home_dist) (home_sqlite) (home_im), inner sep=7pt,
    label={[font=\small\bfseries, text=fuxiBlue, yshift=1pt]above:{Home 节点}}] (home_box) {};
  \node[externalBoundary, fit=(worker1_cc) (worker1_sandbox) (worker1_relay), inner sep=7pt,
    label={[font=\small\bfseries, text=fuxiBlue, yshift=1pt]above:{Worker A（笔记本）}}] (worker1_box) {};
  \node[externalBoundary, fit=(worker2_cc) (worker2_relay), inner sep=7pt,
    label={[font=\small\bfseries, text=fuxiBlue, yshift=-1pt]below:{Worker B（家用台式机）}}] (worker2_box) {};
\end{scope}

% ── Home 节点内部连线 ──
\draw[flow] (home_xuannv) -- (home_dist);
\draw[flow] (home_dist) -- (home_sqlite);
\draw[flow] (home_im) -- (home_dist);

% ── Worker 内部连线 ──
\draw[flow] (worker1_cc) -- (worker1_sandbox);
\draw[flow] (worker1_sandbox) -- (worker1_relay);
\draw[flow] (worker2_cc) -- (worker2_relay);

% ── 跨节点通道（粗虚线 + HMAC 标签）──
% 编号气泡放在 home_dist 右侧紧贴出口，远离 worker box
\node[numBubble] at ($(home_dist.east)+(4mm,2mm)$) {1};
\draw[flowEvent, very thick, dash pattern=on 3pt off 2pt]
  (home_dist.east) -- node[above=2pt, font=\scriptsize, color=fuxiBlue, sloped, midway] {HMAC + ssh tunnel}
  (worker1_relay.west);

\node[numBubble] at ($(home_dist.east)+(4mm,-2mm)$) {2};
\draw[flowEvent, very thick, dash pattern=on 3pt off 2pt]
  (home_dist.east) to[bend right=15] node[below=2pt, font=\scriptsize, color=fuxiBlue, sloped, pos=0.4] {/dist/* HMAC}
  (worker2_relay.west);

% ── auto-pin 注释（右下角小框）──
\node[anchor=north west, font=\scriptsize\itshape, text=fuxiGrayDark, align=left,
      draw=fuxiGrayLine!40, fill=fuxiGrayBox!30, inner sep=5pt, line width=0.3pt]
  at ($(worker2_box.south west)+(0,-9mm)$) {
    auto-pin: $\arg\min_{n\in\mathcal{N}_p}\dfrac{n.\,\mathrm{inflight}}{n.\,\mathrm{max\_concurrency}}$\\
    用户显式 \texttt{@<node>} 优先 / 守 \texttt{task.pinned\_node}
  };

\end{tikzpicture}
\end{document}
